
% 附录源 PDF 第 11—17 页；A.3.3 全部证明。
\subsection{定理 18.1 的证明}

为简化记号，令 $n=(n-n_l+1)$，并将关于 $t_0$ 的求和范围记为从 1 到 $n$。有
\begin{align*}
0\le{}&d_{0n}(\widehat\Psi_{k_n},\psi_0)\\
={}&\frac1n\sum_{t_0}P_{\theta_{0,t_0},Z(t_0)}
 \{L(\widehat\Psi_{k_n}(P_{t_0-1}))-L(\psi_0)\}\\
&-(1+\delta)\frac1n\sum_{t_0}
 \{L(\widehat\Psi_{k_n}(P_{t_0-1}))-L(\psi_0)\}(O(t_0),Z(t_0))\\
&+(1+\delta)\frac1n\sum_{t_0}
 \{L(\widehat\Psi_{k_n}(P_{t_0-1}))-L(\psi_0)\}(O(t_0),Z(t_0))\\
\le{}&\frac1n\sum_{t_0}P_{\theta_{0,t_0},Z(t_0)}
 \{L(\widehat\Psi_{k_n}(P_{t_0-1}))-L(\psi_0)\}\\
&-(1+\delta)\frac1n\sum_{t_0}
 \{L(\widehat\Psi_{k_n}(P_{t_0-1}))-L(\psi_0)\}(O(t_0),Z(t_0))\\
&+(1+\delta)\frac1n\sum_{t_0}
 \{L(\widehat\Psi_{\widetilde k_n}(P_{t_0-1}))-L(\psi_0)\}(O(t_0),Z(t_0))\\
={}&\frac1n\sum_{t_0}P_{\theta_{0,t_0},Z(t_0)}
 \{L(\widehat\Psi_{k_n}(P_{t_0-1}))-L(\psi_0)\}\\
&-(1+\delta)\frac1n\sum_{t_0}
 \{L(\widehat\Psi_{k_n}(P_{t_0-1}))-L(\psi_0)\}\\
&+(1+\delta)\frac1n\sum_{t_0}
 \{L(\widehat\Psi_{\widetilde k_n}(P_{t_0-1}))-L(\psi_0)\}\\
&-(1+2\delta)\frac1n\sum_{t_0}P_{\theta_{0,t_0},Z(t_0)}
 \{L(\widehat\Psi_{\widetilde k_n}(P_{t_0-1}))-L(\psi_0)\}\\
&+(1+2\delta)\frac1n\sum_{t_0}P_{\theta_{0,t_0},Z(t_0)}
 \{L(\widehat\Psi_{\widetilde k_n}(P_{t_0-1}))-L(\psi_0)\}.
\end{align*}
将最后一个表达式的前两项之和记为 $R_{n,k_n}$，第三、四项之和记为 $T_{n,\widetilde k_n}$；最后一项是基准项 $(1+2\delta)d_{0n}(\widehat\Psi_{\widetilde k_n},\psi_0)$。因此，
\begin{equation}
\begin{aligned}
0\le d_{0n}(\widehat\Psi_{k_n},\psi_0)
\le{}&(1+2\delta)d_{0n}(\widehat\Psi_{\widetilde k_n},\psi_0)\\
&+R_{n,k_n}+T_{n,\widetilde k_n}.
\end{aligned}
\tag{A.7}
\end{equation}

\paragraph{将 $R_{n,k}$（以及 $T_{n,k}$）改写为鞅。}
为简化记号，对有关随机变量引入如下记法：
\[
\begin{aligned}
\widetilde H_k&\equiv\frac1n\sum_{t_0}P_{\theta_{0,t_0},Z(t_0)}
 \{L(\widehat\Psi_k(P_{t_0-1}))-L(\psi_0)\},\\
\bar H_k&\equiv\frac1n\sum_{t_0}
 \{L(\widehat\Psi_k(P_{t_0-1}))-L(\psi_0)\}(O(t_0),Z(t_0)).
\end{aligned}
\]
由 $\psi_0$ 的定义，对所有 $k$ 都有 $\widetilde H_k\ge0$。将 $R_{n,k}$ 和 $T_{n,k}$ 改写为
\[
R_{n,k}=(1+\delta)[\widetilde H_k-\bar H_k]-\delta\widetilde H_k,
\]
以及
\[
T_{n,k}=(1+\delta)[\bar H_k-\widetilde H_k]-\delta\widetilde H_k.
\]

\paragraph{忽略可忽略余项后，用具有确定性负平移的鞅和近似 $R_{n,k}$（以及 $T_{n,k}$）。}
具有负平移的鞅和具有良好的指数尾部行为。为了利用这一性质，需要将随机平移量 $\delta\widetilde H_k\ge0$ 换为确定性平移量，并保证该平移量大于一个可控制的常数。下面利用假设 A3 实现这一点。给定 $K$，定义
\[
\widetilde R_{n,k}\equiv\frac{2K^2}{n}\sum_{t_0=1}^n
 P_{\theta_{0,t_0},Z(t_0)}
 \phi\left(\frac{|L^0(\widehat\Psi_k(P_{t_0-1}))|}{K}\right),
\]
其中 $\phi(x)\equiv\exp(x)-x-1$。由引理 A.2，
\[
\widetilde R_{n,k}\le
 M_2\bigl(1/2K^2+1/6M_1K\exp(M_1/K)\bigr)\widetilde H_k.
\]
将该常数记为 $C_1(M_1,M_2,K)$，即
$\widetilde R_{n,k}\le C_1(M_1,M_2,K)\widetilde H_k$。
定义事件
\[
E_{nk}=\{M_{3n}^{-1}<\widetilde H_k/E_0\widetilde H_k<M_{3n}\},
\]
并用 $I_{E_{nk}}$ 表示其指示变量。由假设 A3，当 $n\to\infty$ 时，有
$P_0^n(I_{E_{n,k_n}}=1)\to1$ 和
$P_0^n(I_{E_{n,\widetilde k_n}}=1)\to1$。这还意味着
$P_0^n(\widetilde R_{n,k_n}/E_0\widetilde H_{k_n}<C_1M_{3n})\to1$。
为简化记号，将 $M_{3n}$ 重新定义为 $\max(C_1,1)M_{3n}$。把 $R_{n,k}$ 分解如下：
\[
\begin{aligned}
R_{n,k}={}&(1+\delta)[\widetilde H_k-\bar H_k]I_{E_{n,k}}
 +(1+\delta)[\widetilde H_k-\bar H_k]I_{E_{n,k}^c}\\
&-\delta\widetilde H_k I(\widetilde H_k>M_{3n}^{-1}E_0\widetilde H_k)\\
&-\delta\widetilde H_k I(\widetilde H_k<M_{3n}^{-1}E_0\widetilde H_k)\\
={}&R^*_{n,k}+e_{n,k},
\end{aligned}
\]
其中
\[
\begin{aligned}
R^*_{n,k}\equiv{}&(1+\delta)[\widetilde H_k-\bar H_k]I_{E_{n,k}}\\
&-\delta\widetilde H_k I(\widetilde H_k>M_{3n}^{-1}E_0\widetilde H_k),\\
e_{n,k}\equiv{}&(1+\delta)[\widetilde H_k-\bar H_k]I^c_{E_{n,k}}\\
&-\delta\widetilde H_k I(\widetilde H_k<M_{3n}^{-1}E_0\widetilde H_k).
\end{aligned}
\]
因此 $R_{n,k_n}=R^*_{n,k_n}+e_{n,k_n}$。由假设 A3，当 $n\to\infty$ 时，
$P_0^n(|e_{n,k_n}|=0)\to1$。类似地，
\[
T_{n,k}=T^*_{n,k}+f_{n,k},
\]
其中
\[
\begin{aligned}
T^*_{n,k}\equiv{}&(1+\delta)[\bar H_k-\widetilde H_k]I_{E_{n,k}}\\
&-\delta\widetilde H_k I(\widetilde H_k>M_{3n}^{-1}E_0\widetilde H_k),\\
f_{n,k}\equiv{}&(1+\delta)[\bar H_k-\widetilde H_k]I_{E_{n,k}}\\
&-\delta\widetilde H_k I(\widetilde H_k<M_{3n}^{-1}E_0\widetilde H_k).
\end{aligned}
\]
使用与 $e_{n,k_n}$ 相同的论证，得
$P_0^n(|f_{n,\widetilde k_n}|=0)\to1$。因此
$T_{n,\widetilde k_n}=T^*_{n,\widetilde k_n}+f_{n,\widetilde k_n}$，其中
$f_{n,\widetilde k_n}$ 以趋于 1 的概率等于零。

\paragraph{已得到的结果与尚需完成的步骤。}
令 $Z_{n2}=e_{n,k_n}+f_{n,\widetilde k_n}$，并令
$Z_{n1}=R^*_{n,k_n}+T^*_{n,\widetilde k_n}$。我们已证明
\[
d_{0n}(\widehat\Psi_{k_n},\psi_0)
 \le(1+2\delta)d_{0n}(\widehat\Psi_{\widetilde k_n},\psi_0)+Z_{n1}+Z_{n2},
\]
其中，当 $n\to\infty$ 时，$P_0^n(Z_{n2}=0)\to1$。下面将证明，对于某个明确的有限常数
$C(M_1,M_2,M_{3n},\delta)$，有
\[
EZ_{n1}=ER^*_{n,k_n}+ET^*_{n,\widetilde k_n}
 \le C(M_1,M_2,M_{3n},\delta)\frac{1+\log K(n)}n,
\]
由此即可完成证明。

\paragraph{对 $R^*_{n,k_n}$ 的尾概率求界。}
\subparagraph{步骤 1：得到确定性的负平移量。}
另定义事件 $E_{n,k,1}=\{\widetilde H_k>M_{3n}^{-1}E_0\widetilde H_k\}$。
取 $s>0$，有
\[
\begin{aligned}
&P_0^n(R^*_{n,k_n}>s)\\
&=P_0^n\left(I_{E_{n,k_n}}\{\widetilde H_{k_n}-\bar H_{k_n}\}
 >\frac1{1+\delta}\{s+\delta\widetilde H_{k_n}I_{E_{n,k_n,1}}\}\right)\\
&\le P_0^n\left(I_{E_{n,k_n}}\{\widetilde H_{k_n}-\bar H_{k_n}\}
 >\frac1{1+\delta}\left\{s+\delta M_{3n}^{-1}
 E_0\widetilde H_k\big|_{k=k_n}I_{E_{n,k_n,1}}\right\}\right).
\end{aligned}
\]
这里利用了事件 $E_{n,k_n,1}$ 蕴含
$\widetilde H_{k_n}\ge M_{3n}^{-1}E_0\widetilde H_k|_{k=k_n}$，因此可以用仅通过 $k_n$ 保留随机性的上述界代替随机量 $\widetilde H_{k_n}$。
将最后一个概率中的事件记为 $A_n$，于是最后的界为 $P_0^n(A_n)$。可写成
\[
P_0^n(A_n)=P_0^n(A_n\text{ 且 }I_{E_{n,k_n,1}}=1)
 +P_0^n(A_n\text{ 且 }I_{E_{n,k_n,1}}=0).
\]
注意，当 $I_{E_{n,k_n,1}}=0$ 时，上述事件不等式的右端等于
$s/(1+\delta)>0$，而左端等于 0。因此
$P_0^n(A_n\text{ 且 }I_{E_{n,k_n,1}}=0)=0$。
由此得到 $P_0^n(R^*_{n,k_n}>s)$ 的界：
\[
\begin{aligned}
&P_0^n(R^*_{n,k_n}>s)\\
&\le P_0^n\Biggl(I_{E_{n,k_n}}\{\widetilde H_{k_n}-\bar H_{k_n}\}
 >\frac1{1+\delta}\left\{s+\delta M_{3n}^{-1}E_0\widetilde H_k\big|_{k=k_n}\right\},\\
&\hspace{22em}\text{且 }E_{n,k_n,1}=1\Biggr)\\
&\le P_0^n\left(I_{E_{n,k_n}}\{\widetilde H_{k_n}-\bar H_{k_n}\}
 >\frac1{1+\delta}\left\{s+\delta M_{3n}^{-1}E_0\widetilde H_k\big|_{k=k_n}\right\}\right)\\
&\le K(n)\max_kP_0^n\left(I_{E_{n,k}}\{\widetilde H_k-\bar H_k\}
 >\frac1{1+\delta}\{s+\delta M_{3n}^{-1}E_0\widetilde H_k\}\right).
\end{aligned}
\]
最后一个不等式利用了如下事实：对于一组随机变量 $(X(k):k)$、常数 $(c(k):k)$ 和随机指标 $k_n$，有
\[
\begin{aligned}
P_0^n(X(k_n)<c(k_n))
&\le P_0^n(\text{至少存在一个 }k\text{ 使 }X(k)<c(k))\\
&\le\sum_{k=1}^{K(n)}P_0^n(X(k)<c(k))\\
&\le K(n)\max_kP_0^n(X(k)<c(k)).
\end{aligned}
\]
类似地，对 $T^*_{n,\widetilde k_n}$，有
\[
\begin{aligned}
&P_0^n(T^*_{n,\widetilde k_n}>s)\\
&\le K(n)\max_kP_0^n\left(I_{E_{n,k}}\{\bar H_k-\widetilde H_k\}
 >\frac1{1+\delta}\{s+\delta M_{3n}^{-1}E_0\widetilde H_k\}\right).
\end{aligned}
\]

\subparagraph{步骤 2：应用鞅尾概率指数不等式。}
有 $\bar H_k-\widetilde H_k$ 等于鞅和
$n^{-1}\sum_{t_0}Z_{k,t_0}-E(Z_{k,t_0}\mid\bar O(t_0-1))$，其中
\[
Z_{k,t_0}=\{L(\widehat\Psi_k(P_{t_0-1}))-L(\psi_0)\}(\bar O(t_0)).
\]
由假设 A1，随机变量 $Z_{k,t_0}$ 有界：$|Z_{k,t_0}|\le M_1$ 几乎处处成立。
现在，对每个 $k$ 分别取 $R=R_k=M_{3n}E_0\widetilde H_k$，将定理 A.2 的鞅不等式，具体而言是式 (A.6)，应用于 $\widetilde H_k-\bar H_k$。由于这样选择 $R$，在 $s>0$ 且 $s$ 较小时，所得尾概率呈现 $\exp(-cM_{3n}ns)$ 的行为，而通常为 $\exp(-cns^2)$。这进而表明，余项 $R^*_{n,k_n}$ 和 $T^*_{n,\widetilde k_n}$ 的期望按 $\log K(n)/n$ 的速率收敛，而通常的速率为 $\log K(n)/\sqrt n$。
为便于查阅，下面把给定 $k$ 时的鞅指数不等式明确列出。

\begin{bookresult}{引理}{A.3}
固定 $K$，令
\[
\widetilde R_{n,k}\equiv\frac{2K^2}{n}\sum_{i=1}^nP_{\theta_{0,i},Z(i)}
 \phi\left(\frac{|L^0(\widehat\Psi_k(P_{i-1}))|}{K}\right),
\]
其中 $\phi(x)\equiv\exp(x)-x-1$。令
\[
\begin{aligned}
M_{n,k}=\sum_{i=1}^n\Bigl\{
&L^0(\widehat\Psi_k(P_{i-1}))(Z(i),O(i))\\
&-E_0(L^0(\widehat\Psi_k(P_{i-1}))\mid Z(i))\Bigr\}.
\end{aligned}
\]
对于任意满足 $c_0^2\ge C^2(c_1+1)$ 的 $R_k,c_0,c_1\ge0$，以及
$\alpha\in(c_0\sqrt{R_k}/\sqrt n,c_1R_k/K)$，有
\[
P\left(I(\widetilde R_{n,k}\le R_k)\frac{M_{n,k}}n\ge\alpha\right)
 \le2\exp\left\{-\frac{n\alpha^2}{C^2(K\alpha+R_k)}\right\}.
\]
\end{bookresult}

为了在给定
\[
\alpha(s)=\frac1{1+\delta}(s+\delta M_{3n}^{-1}E_0\widetilde H_k)
\]
处，将这一不等式应用于 $I(\widetilde R_{n,k}<R_k)(\bar H_k-\widetilde H_k)$ 的上述尾概率，其中 $R_k=M_{3n}E_0\widetilde H_k$，需要能够选择满足 $c_0^2\ge C^2(c_1+1)$ 的 $c_0,c_1$，使得
\begin{equation}
\frac{c_0\sqrt{R_k}}{\sqrt n}
 \le\frac1{1+\delta}[s+\delta M_{3n}^{-1}E_0\widetilde H_k]
 \le c_1\frac{R_k}{K}.
\tag{A.8}
\end{equation}
注意，$M_{3n}^{-1}E_0\widetilde H_k=M_{3n}^{-2}R_k$，因此需要在
\[
\alpha(s)=1/(1+\delta)(s+\delta M_{3n}^{-2}R_k)
\]
处应用该不等式。
所以，要选择 $c_0,c_1$，使得
$\alpha(s)\in(c_0R_k^{0.5}/n^{0.5},c_1R_k/K)$。取
$c_0^2=c_0^2(c_1)=C^2(c_1+1)$。由于鞅过程
$\bar H_k-\widetilde H_k$ 被 $2M_1$ 界定，当 $c_1R_k/K>2M_1$ 时，上界就不再构成限制。因此选择
$c_1=c_1(M_1)=2M_1K/R_k$，从而保证对所有 $s$，$\alpha(s)$ 均无上界限制。
令 $c_0(M_1)=c_0(c_1(M_1))$ 为相应的 $c_0$。于是，对于任意
$\alpha\in(c_0(M_1)R_k^{0.5}n^{-0.5},\infty)$，都有
\[
P_0^n\bigl(I(\widetilde R_{nk}<R_k)(\bar H_k-\widetilde H_k)>\alpha\bigr)
 \le2\exp\{-n\alpha^2/[C^2(K\alpha+R_k)]\}.
\]
左侧不等式 $\alpha(s)>c_0(M_1)R_k^{0.5}n^{-0.5}$ 等价于
\begin{equation}
s>-\delta M_{3n}^{-2}R_k
 +C^2(c_1(M_1)+1)n^{-0.5}R_k^{0.5}(1+\delta).
\tag{A.9}
\end{equation}
右侧第一项为负，以 $M_{3n}^{-2}R_k$ 的速率趋于零；第二项为正，以 $R_k^{0.5}n^{-0.5}$ 的速率趋于零。由假设 A4，
\[
\max_k\frac{R_k^{0.5}n^{-0.5}}{M_{3n}^{-2}R_k}\longrightarrow0.
\]
因此，$n$ 足够大时，式 (A.9) 的右端为负。这就证明了
$\alpha(s)>c_0(M_1)R_k^{0.5}n^{-0.5}$ 对所有 $s>0$ 成立。
于是存在 $n_1$，使得对所有 $n>n_1$ 及所有 $s>0$，
\[
\begin{aligned}
&P_0^n\bigl(I(\widetilde R_{nk}<R_k)(\bar H_k-\widetilde H_k)>\alpha(s)\bigr)\\
&\le2\exp\{-n\alpha(s)^2/[C^2(K\alpha(s)+R_k)]\}\\
&=2\exp\left(-C^{-2}\frac n{(1+\delta)^2}
 \frac{(s+\delta M_{3n}^{-2}R_k)^2}
 {R_k+\frac K{1+\delta}(s+\delta M_{3n}^{-2}R_k)}\right).
\end{aligned}
\]

\subparagraph{步骤 3：考察尾概率的渐近行为。}
注意，
\[
\begin{aligned}
&\frac{(s+\delta M_{3n}^{-2}R_k)^2}
 {R_k+\frac K{1+\delta}(s+\delta M_{3n}^{-2}R_k)}\\
&=\frac{s+\delta M_{3n}^{-2}R_k}
 {\frac{R_k}{s+\delta M_{3n}^{-2}R_k}+\frac K{1+\delta}}\\
&\ge\frac{s+\delta M_{3n}^{-2}R_k}{M_{3n}^2/\delta+K}
 \ge\frac s{M_{3n}^2/\delta+K}\\
&=M_{3n}^{-2}\frac s{\delta^{-1}+M_{3n}^{-2}K}
 \ge M_{3n}^{-2}\frac s{\delta^{-1}+K}.
\end{aligned}
\]
这里利用了对所有 $n$ 都有 $M_{3n}>1$，因而 $KM_{3n}^{-2}\le K$。这说明，对于 $s>0$，
\[
P_0^n(R^*_{n,k_n}>s)\le2K(n)
 \exp\left(-\frac{nM_{3n}^{-2}}{c(M_1,M_2,\delta)}s\right),
\]
其中 $c(M_1,M_2,\delta)=2C^2(1+\delta)^2(K+\delta^{-1})$。

\paragraph{根据尾概率界，对 $R^*_{n,k_n}$ 的期望求界。}
由于
$ER^*_{n,k_n}\le\int_0^\infty P_0^n(R^*_{n,k_n}>s)\,\mathrm ds$，对每个 $u>0$，都有
\[
ER^*_{n,k_n}\le u+\int_u^\infty2K(n)
 \exp\left(-\frac{M_{3n}^{-2}n}{c(M_1,M_2,\delta)}s\right)\,\mathrm ds.
\]
该上界在
\[
u_n=\frac{c(M_1,M_2,\delta)\log(2K(n))}{nM_{3n}^{-2}}
\]
处取到最小值，最小值为
\[
\frac{c(M_1,M_2,\delta)(\log(2K(n))+1)}{nM_{3n}^{-2}}.
\]
因此，
\[
ER^*_{n,k_n}\le c(M_1,M_2,\delta)
 \frac{1+\log(2K(n))}{nM_{3n}^{-2}}.
\]
类似地，可以得到 $ET_{n,\widetilde k_n}$ 的这一界。这就证明了假设 A4 成立时的定理。

\paragraph{假设 A4 不成立时的情形。}
现在考虑假设 A4 不成立的情况。此时主导项
$E_0d_{0n}(\widehat\Psi_{\widetilde k_n},\psi_0)=O(n^{-1}M_{3n}^3)$。
首先考虑式 (A.9) 右端为负的情况。此时，对所有 $s>0$，前面已经得到了所需的 $P_0^n(R^*_{n,k}>s)$ 界。
再考虑式 (A.9) 右端为正的情况。此时，
\[
R_k^{0.5}<(1+\delta)\delta^{-1}C^2(c_1(M_1)+1)M_{3n}^2n^{-0.5},
\]
从而式 (A.9) 的右端可以由下式界定：
\[
cM_{3n}^2n^{-1}\equiv
 (1+\delta)\delta^{-1}C^4(c_1(M_1)+1)^2M_{3n}^2n^{-1}.
\]
因此，在这一情形下，对这个明确的常数 $c>0$，当
$s>cM_{3n}^2n^{-1}$ 时，即有
$P_0^n(I(\widetilde R_{nk}<R_k)(\bar H_k-\widetilde H_k)>\alpha(s))$
所需的指数界。
接下来，对任意 $u>cM_{3n}^2n^{-1}$，有
\[
\begin{aligned}
E_0R^*_{n,k_n}
={}&\int_0^uP_0^n(R^*_{n,k_n}>s)\,\mathrm ds
 +\int_u^\infty P_0^n(R^*_{n,k_n}>s)\,\mathrm ds\\
\le{}&u+K(n)\max_k\int_u^\infty P_0^n(R^*_{n,k}>s)\,\mathrm ds\\
\le{}&u+2K(n)\int_u^\infty
 \exp\left(-\frac{nM_{3n}^{-2}}{c(M_1,M_2,\delta)}s\right)\,\mathrm ds\\
={}&u+2K(n)\frac{c(M_1,M_2,\delta)}{nM_{3n}^{-2}}
 \exp\left(-\frac{nM_{3n}^{-2}}{c(M_1,M_2,\delta)}u\right).
\end{aligned}
\]
最优的 $u$ 为
\[
u^*=\max\bigl(cM_{3n}^2n^{-1},
 c(M_1,M_2,\delta)M_{3n}^2n^{-1}\log(2K(n))\bigr).
\]
若 $c>c(M_1,M_2,\delta)\{\log2+\log K(n)\}$，则
$u^*=cM_{3n}^2n^{-1}$。注意，这还意味着
$K(n)<\exp(cc(M_1,M_2,\delta)^{-1})$。将该 $u^*$ 代入最后的表达式，得到第一项为 $u^*$，再加上
\[
2K(n)c(M_1,M_2,\delta)^{-1}n^{-1}M_{3n}^2
 \exp(-c/c(M_1,M_2,\delta)).
\]
利用 $K(n)$ 的界，可知最终表达式为 $O(M_{3n}^2n^{-1})$。
再假设 $c<c(M_1,M_2,\delta)\{\log2+\log K(n)\}$，则
$u^*=c(M_1,M_2,\delta)M_{3n}^2n^{-1}(\log2+\log K(n))$。
代入此时的 $u^*$，可知最终表达式为
$O(M_{3n}^2n^{-1}(1+\log K(n)))$。
因此，无论哪一种情形，都有
$E_0R^*_{n,k_n}<C_1M_{3n}^2n^{-1}(1+\log K(n))$，其中
$C_1=C_1(M_1,M_2,\delta)<\infty$ 为某个通用常数。
对 $E_0T^*_{n,\widetilde k_n}$ 可以采用相同的求界方法。
由此证明了，在假设 A4 不成立时，
\[
\begin{aligned}
d_{0n}(\widehat\Psi_{k_n},\psi_0)
={}&o_P(n^{-1}M_{3n}^3)\\
&+o_P\bigl(n^{-1}M_{3n}^2(1+\log K(n))\bigr).
\end{aligned}
\]
至此，定理 18.1 得证。
